\documentclass[11pt,reqno]{amsart}
\usepackage[utf8]{inputenc}
\usepackage[T1]{fontenc}
\usepackage{amsmath,amssymb,amsthm,mathtools}
\usepackage[margin=2.6cm]{geometry}
\usepackage{hyperref}

\theoremstyle{plain}
\newtheorem{theorem}{Theorem}[section]
\newtheorem{proposition}[theorem]{Proposition}
\newtheorem{lemma}[theorem]{Lemma}
\newtheorem{corollary}[theorem]{Corollary}
\theoremstyle{definition}
\newtheorem{definition}[theorem]{Definition}
\newtheorem{remark}[theorem]{Remark}
\newtheorem{heuristic}[theorem]{Heuristic}

\DeclareMathOperator{\sgn}{sgn}
\newcommand{\Ms}{M_{\mathrm{s}}}
\newcommand{\Ap}{\mathcal{A}_{+}}
\newcommand{\Am}{\mathcal{A}_{-}}
\newcommand{\Anet}{\mathcal{A}}
\newcommand{\eps}{\varepsilon}

\begin{document}

\title{The area function of the Mertens excursion skeleton}
\author{}
\date{}
\maketitle

\section{The area function of the excursion skeleton}
\label{sec:excursion-area}

\noindent\textbf{Convention on RH.} Every statement carries one of three
labels: \emph{[unconditional]}, \emph{[under RH]} (the Riemann hypothesis is
assumed; simple zeros are stated explicitly when also needed), or
\emph{[heuristic]} (not proved; a model). The definitions and the exact
identities of \S\ref{subsec:def}--\S\ref{subsec:formulas} are
unconditional: they use no hypothesis on the zeros of $\zeta$.

\subsection{Definitions}
\label{subsec:def}

Let $\mu$ be the Möbius function and
\[
  M(x)=\sum_{n\le x}\mu(n),\qquad x\ge 0,
\]
the Mertens function, piecewise constant with jumps at the integers. Let
$0=z_0<z_1<z_2<\dots$ be the points where $M$ vanishes. Since the jumps of
$M$ are $0$ or $\pm1$, every sign change passes through a zero, and $M$
changes sign infinitely often \emph{[unconditional]}. The $\nu$-th
\emph{excursion} is the interval
\[
  I_\nu=(u_\nu,v_\nu],\qquad u_\nu=z_{\nu-1},\quad v_\nu=z_\nu,
  \qquad \ell_\nu=v_\nu-u_\nu ,
\]
in whose interior $M$ does not vanish. Inside $I_\nu$ we set
\[
  H_\nu=\max_{u_\nu<n\le v_\nu}|M(n)|,\qquad
  \tau_\nu=\min\{n\in I_\nu:\ |M(n)|=H_\nu\},\qquad
  \sigma_\nu=\sgn M(\tau_\nu)\in\{\pm1\},
\]
and we write $a_\nu=\tau_\nu-u_\nu$, $b_\nu=v_\nu-\tau_\nu$, so that
$\ell_\nu=a_\nu+b_\nu$. Since $M$ needs at least $H_\nu$ steps to climb from
$0$ to $H_\nu$ and as many to return, one has
$H_\nu\le\min(a_\nu,b_\nu)$ \emph{[unconditional]}.

\begin{definition}[Regularized skeleton]
\label{def:Ms}
The \emph{skeleton} $\Ms$ is the continuous, piecewise linear function
defined on each excursion by
\[
  \Ms(t)=
  \begin{cases}
    \sigma_\nu\sqrt{H_\nu}\,\dfrac{t-u_\nu}{a_\nu}, & u_\nu\le t\le \tau_\nu,\\[2ex]
    \sigma_\nu\sqrt{H_\nu}\,\dfrac{v_\nu-t}{b_\nu}, & \tau_\nu\le t\le v_\nu .
  \end{cases}
\]
On the terminal block $(u,N]$, if $M(N)\neq0$, the descent ends at
$\sigma\sqrt{|M(N)|}$ instead of at $0$. We also write $L$ for the
\emph{unregularized} skeleton, defined in the same way but with peak
$\sigma_\nu H_\nu$ instead of $\sigma_\nu\sqrt{H_\nu}$.
\end{definition}

\begin{definition}[The three curves]
\label{def:areas}
\[
  \Ap(x)=\int_0^x \Ms^{+}(t)\,dt,\qquad
  \Am(x)=-\int_0^x \Ms^{-}(t)\,dt,\qquad
  \Anet(x)=\Ap(x)+\Am(x)=\int_0^x \Ms(t)\,dt,
\]
where $\Ms^{\pm}=\max(\pm \Ms,0)$. In the figure, $\Ap$ is the blue curve,
$\Am$ the red one and $\Anet$ the black one. $\Ap$ is nondecreasing, $\Am$ is
nonincreasing, and both are constant on the excursions of the opposite sign:
this is what produces the plateaus.
\end{definition}

\subsection{Exact formulas}
\label{subsec:formulas}

\begin{lemma}[Increment per excursion]
\label{lem:jump}
\emph{[unconditional]} For every complete excursion,
\[
  \Anet(v_\nu)-\Anet(u_\nu)=\sigma_\nu\,\frac{\sqrt{H_\nu}\,\ell_\nu}{2}.
\]
The discrete version agrees: $\sum_{n\in I_\nu}\Ms(n)
=\sigma_\nu\sqrt{H_\nu}\bigl((a_\nu+1)+(b_\nu-1)\bigr)/2
=\sigma_\nu\sqrt{H_\nu}\,\ell_\nu/2$.
\end{lemma}

\begin{proof}
The two pieces of $\Ms$ are triangles of height $\sqrt{H_\nu}$ with bases
$a_\nu$ and $b_\nu$; their areas are $\sqrt{H_\nu}a_\nu/2$ and
$\sqrt{H_\nu}b_\nu/2$. In the discrete sum, the ascent contributes
$\sqrt{H_\nu}\sum_{j=1}^{a_\nu}j/a_\nu=\sqrt{H_\nu}(a_\nu+1)/2$ and the
descent $\sqrt{H_\nu}\sum_{j=0}^{b_\nu-1}j/b_\nu=\sqrt{H_\nu}(b_\nu-1)/2$.
\end{proof}

\begin{proposition}[Explicit form of the three curves]
\label{prop:explicit}
\emph{[unconditional]} Let $x\in I_\nu$. Then
\[
  \Anet(x)=
  \begin{cases}
    \Anet(u_\nu)+\sigma_\nu\sqrt{H_\nu}\dfrac{(x-u_\nu)^2}{2a_\nu},
      & u_\nu\le x\le\tau_\nu,\\[2.4ex]
    \Anet(u_\nu)+\sigma_\nu\sqrt{H_\nu}
      \left(\dfrac{a_\nu}{2}+(x-\tau_\nu)-\dfrac{(x-\tau_\nu)^2}{2b_\nu}\right),
      & \tau_\nu\le x\le v_\nu,
  \end{cases}
\]
and, at the zeros of $M$,
\[
  \Anet(v_k)=\sum_{\nu\le k}\sigma_\nu\frac{\sqrt{H_\nu}\,\ell_\nu}{2},
  \qquad
  \Ap(v_k)=\sum_{\substack{\nu\le k\\ \sigma_\nu=+1}}\frac{\sqrt{H_\nu}\,\ell_\nu}{2},
  \qquad
  \Am(v_k)=-\sum_{\substack{\nu\le k\\ \sigma_\nu=-1}}\frac{\sqrt{H_\nu}\,\ell_\nu}{2}.
\]
Thus the three curves are piecewise quadratic, of class $C^1$, and are
determined by the data $(u_\nu,\tau_\nu,v_\nu,H_\nu,\sigma_\nu)$.
\end{proposition}

\begin{lemma}[Critical points and inflections]
\label{lem:critical}
\emph{[unconditional]} One has $\Anet'=\Ms$, so the local extrema of
$\Anet$ are exactly the zeros of $M$, and $\Anet$ is monotone on each
excursion. Since $\Ms$ is piecewise linear, $\Anet''$ is piecewise constant
and changes sign at $\tau_\nu$: the inflection points of $\Anet$ are the
first argmaxima of $|M|$.
\end{lemma}



\subsection{Internal correlations of an excursion}
\label{subsec:internal-corr}

The excursion decomposition also reorganizes the quadratic form of the
Mertens function.  This is the local correlation structure that will be used
below.  Consider first an arbitrary finite block $x_1,\dots,x_L\in\{-1,0,1\}$
and its partial sums
\[
  S_t=\sum_{j=1}^t x_j,\qquad 1\le t\le L,
\]
and define the internal correlations
\[
  C_h=\sum_{j=1}^{L-h}x_jx_{j+h},\qquad 1\le h<L.
\]

\begin{lemma}[Triangular correlation identity]
\label{lem:triangular-corr}
\emph{[unconditional]} For every finite block,
\begin{equation}
\label{eq:general-triangular}
  \sum_{t=1}^{L}S_t^2
  =S_L^2+S_L\sum_{t=1}^{L-1}S_t-\sum_{h=1}^{L-1}hC_h.
\end{equation}
In particular, for a bridge $S_L=0$,
\begin{equation}
\label{eq:bridge-triangular}
  \sum_{t=1}^{L-1}S_t^2=-\sum_{h=1}^{L-1}hC_h.
\end{equation}
\end{lemma}

\begin{proof}
One has
\[
  \sum_{t=1}^{L}S_t^2
  =\sum_{i,j=1}^{L}x_ix_j\bigl(L+1-\max(i,j)\bigr).
\]
Using $\max(i,j)=\tfrac12(i+j+|i-j|)$, the $|i-j|$ term gives
$-\sum_{h<L}hC_h$, whereas the $i+j$ term is reduced using
\[
  \sum_{i=1}^{L} i x_i
  =LS_L-\sum_{t=1}^{L-1}S_t.
\]
Collecting terms yields \eqref{eq:general-triangular}; the bridge case is
immediate.
\end{proof}

For a complete excursion $I_\nu=(u_\nu,v_\nu]$, put $L=\ell_\nu$ and orient
its increments by
\[
  x_j=\sigma_\nu\mu(u_\nu+j),\qquad
  S_t=\sum_{j=1}^{t}x_j=|M(u_\nu+t)|.
\]
Then $S_t>0$ for $1\le t<L$ and $S_L=0$.  Write
\[
  Q_\nu:=\sum_{t=1}^{\ell_\nu-1}S_t^2,
  \qquad
  C_h^{(\nu)}:=\sum_{j=1}^{\ell_\nu-h}x_jx_{j+h}.
\]
Equation \eqref{eq:bridge-triangular} gives the exact local identity
\begin{equation}
\label{eq:Qcorr}
  Q_\nu=-\sum_{h=1}^{\ell_\nu-1}hC_h^{(\nu)}\ge0.
\end{equation}

\begin{lemma}[Zeroth moment of the internal correlations]
\label{lem:zeroth-corr}
\emph{[unconditional]} If
$q_\nu=\sum_{j=1}^{\ell_\nu}x_j^2$ is the number of non-zero steps in the
excursion, then
\begin{equation}
\label{eq:zeroth-corr}
  \sum_{h=1}^{\ell_\nu-1}C_h^{(\nu)}=-\frac{q_\nu}{2}.
\end{equation}
\end{lemma}

\begin{proof}
Since $\sum_jx_j=0$,
\[
  0=\left(\sum_jx_j\right)^2
   =\sum_jx_j^2+2\sum_{h=1}^{\ell_\nu-1}C_h^{(\nu)}.
\]
\end{proof}

\begin{proposition}[Quadratic decomposition by excursions]
\label{prop:quadratic-decomp}
\emph{[unconditional]} If $x$ is a zero of $M$, then
\begin{equation}
\label{eq:quadratic-decomp-zero}
  \sum_{n\le x}M(n)^2
   =\sum_{v_\nu\le x}Q_\nu
   =-\sum_{v_\nu\le x}\ \sum_{h=1}^{\ell_\nu-1}hC_h^{(\nu)}.
\end{equation}
For general $x$, the same formula holds with the terminal block added, and
its contribution is given by \eqref{eq:general-triangular}.
\end{proposition}

\begin{remark}[Blindness to pairs crossing excursion boundaries]
\label{rem:blindness-quadratic}
Globally,
\[
  \sum_{n\le x}M(n)^2
  =\sum_{a,b\le x}\mu(a)\mu(b)\bigl(x+1-\max(a,b)\bigr),
\]
which contains all pairs $a,b\le x$.  After partitioning at the zeros of
$M$, Proposition~\ref{prop:quadratic-decomp} rewrites the same form using
only pairs lying in the same excursion.  Pairs crossing zeros are therefore
reorganized away by the bridge identities.  This is a quadratic analogue
of an excursion-boundary blindness phenomenon: the triangular form is
completely determined by the internal blocks plus the terminal block.
\end{remark}

Let
\[
  \beta_\nu:=\sum_{t=1}^{\ell_\nu-1}S_t
             =\sum_{n\in I_\nu}|M(n)|.
\]
By Cauchy--Schwarz,
\begin{equation}
\label{eq:beta-Q}
  \beta_\nu^2\le \ell_\nu Q_\nu
  =-\ell_\nu\sum_{h=1}^{\ell_\nu-1}hC_h^{(\nu)}.
\end{equation}
Thus the first weighted moment of the internal correlations controls the
area of the excursion.  The converse need not hold without additional
information on the flatness of the profile.

To separate the support from the sign geometry, define
\[
  P_h^{(\nu)}:=\#\{1\le j\le \ell_\nu-h:\ x_jx_{j+h}\ne0\},
\]
and let $D_h^{(\nu)}$ be the number of those pairs with opposite signs.
Then
\begin{equation}
\label{eq:C-P-D}
  C_h^{(\nu)}=P_h^{(\nu)}-2D_h^{(\nu)},
\end{equation}
and hence
\begin{equation}
\label{eq:Q-P-D}
  Q_\nu=\sum_{h=1}^{\ell_\nu-1}
  h\bigl(2D_h^{(\nu)}-P_h^{(\nu)}\bigr).
\end{equation}
This decomposition isolates the squarefree support $P_h^{(\nu)}$ from the
ordering of signs encoded by $D_h^{(\nu)}$.

\begin{remark}[Height versus quadratic energy]
\label{rem:height-Q}
If an excursion reaches height $H_\nu$, then the least possible value of
$Q_\nu$ is obtained, up to the zero steps, by climbing and descending at
maximal speed.  Consequently
\[
  Q_\nu\ge 2\sum_{r=1}^{H_\nu-1}r^2+H_\nu^2
  =\frac{2}{3}H_\nu^3+O(H_\nu^2).
\]
Therefore a local bound $Q_\nu\ll \ell_\nu^{2+o(1)}$ gives only
$H_\nu\ll \ell_\nu^{2/3+o(1)}$.  The natural primary object is thus the
area normalization $\beta_\nu/\ell_\nu^{3/2}$ rather than the stronger
height normalization $H_\nu/\sqrt{\ell_\nu}$.
\end{remark}

\subsection{Where $\zeta$ enters}
\label{subsec:zeta}

The black curve depends on $M$ nonlinearly, through the square root of the
peak, so $\Anet$ has no Mellin transform algebraic in $\zeta$. Its
dependence on $\zeta$ is that of the data in
Proposition~\ref{prop:explicit}: the nodes $u_\nu,v_\nu$ are the zeros of
$M$, the signs $\sigma_\nu$ are those of $M$ on each excursion, and
$H_\nu=\max_{I_\nu}|M|$. With the explicit formula of
Titchmarsh~\cite[\S14.27]{Titchmarsh}, valid \emph{[under RH and simple
zeros]}, the limit being taken along a suitable sequence $T_k\to\infty$ and
$x$ not an integer,
\begin{equation}
\label{eq:explicitM}
  M(x)=\sum_{\rho}\frac{x^{\rho}}{\rho\,\zeta'(\rho)}-2
   +\sum_{n\ge1}\frac{(-1)^{n-1}(2\pi/x)^{2n}}{n\,(2n)!\,\zeta(2n+1)} ,
\end{equation}
the three ingredients above are written in terms of the zeros of $\zeta$.

For the unregularized skeleton $L$, by contrast, the dependence is linear.
If $\widetilde{A}(x)=\int_0^x M(t)\,dt$, then
$\int_0^\infty \widetilde{A}(x)x^{-s-2}\,dx=\dfrac{1}{s(s+1)\zeta(s)}$ for
$\Re s>1$, and by Mellin inversion, computing the residues at $\rho$, at
$s=0$ (with $\zeta(0)=-\tfrac12$) and at $s=-1$ (with
$\zeta(-1)=-\tfrac1{12}$),
\begin{equation}
\label{eq:explicitA}
  \widetilde{A}(x)=\sum_{\rho}\frac{x^{\rho+1}}{\rho(\rho+1)\zeta'(\rho)}
  -2x+12+O(1/x)
  \qquad\text{\emph{[under RH and simple zeros]}},
\end{equation}
a series that converges better than \eqref{eq:explicitM} thanks to the
factor $|\rho|^{-2}$. Moreover, with $R=M-L$ one has the unconditional
identity
\begin{equation}
\label{eq:dipoles}
  \int_0^{v_k}L
  =\widetilde{A}(v_k)+\sum_{\text{legs}\ \subseteq (0,v_k]}\ \sum_{m\in\text{leg}}\mu(m)\,(m-m_{\text{leg}}),
\end{equation}
where each leg is an ascent or a descent and $m_{\text{leg}}$ is its
midpoint: the area of the residual on each leg is exactly a dipole of $\mu$.

\subsection{Bounds}
\label{subsec:bounds}

\begin{lemma}[Pointwise transfer]
\label{lem:transfer}
\emph{[unconditional]} If $|M(t)|\le c\,t^{1/2}$ for all $t\le x$, then
$|\Ms(t)|\le c^{1/2}t^{1/4}$ for all $t\le x$, and therefore
\[
  \Ap(x)+|\Am(x)|\le \tfrac45\,c^{1/2}\,x^{5/4},
  \qquad |\Anet(x)|\le \tfrac45\,c^{1/2}\,x^{5/4}.
\]
\end{lemma}

\begin{proof}
Let $t\in I_\nu$. If $t\ge\tau_\nu$, then
$|\Ms(t)|\le\sqrt{H_\nu}\le c^{1/2}\tau_\nu^{1/4}\le c^{1/2}t^{1/4}$. If
$t\le\tau_\nu$, then $|\Ms(t)|$ is the line joining $(u_\nu,0)$ to
$(\tau_\nu,\sqrt{H_\nu})$, and $\sqrt{H_\nu}\le c^{1/2}\tau_\nu^{1/4}$;
since $t\mapsto c^{1/2}t^{1/4}$ is concave and nonnegative at $u_\nu$, it
lies above that chord. Integrating,
$\int_0^x c^{1/2}t^{1/4}dt=\tfrac45 c^{1/2}x^{5/4}$.
\end{proof}

\begin{corollary}[Unconditional bound]
\label{cor:uncond}
\emph{[unconditional]} By the Vinogradov--Korobov zero-free region,
$M(x)\ll x\exp\bigl(-c(\log x)^{3/5}(\log\log x)^{-1/5}\bigr)$, whence
\[
  \Ap(x)+|\Am(x)|\ll x^{3/2}\exp\Bigl(-\tfrac{c}{2}(\log x)^{3/5}(\log\log x)^{-1/5}\Bigr).
\]
\end{corollary}

\begin{corollary}[Bound under RH]
\label{cor:rh}
\emph{[under RH]} $M(x)\ll_\eps x^{1/2+\eps}$ implies
$\Ap(x)+|\Am(x)|\ll_\eps x^{5/4+\eps}$. Conversely, if
$\Ap(x)+|\Am(x)|\gg x^{5/4+\eps}$ along a sequence $x\to\infty$ for some
fixed $\eps>0$, then RH is false.
\end{corollary}

\begin{remark}
\label{rem:onedirection}
The converse goes only in that direction. Since $\ell_\nu\ge 2H_\nu$, a
bound $\sqrt{H_\nu}\ell_\nu\ll x^{5/4+\eps}$ yields only
$H_\nu\ll x^{5/6+\eps}$, far from RH. The black curve is a refutation test,
not a criterion.
\end{remark}

\begin{remark}[Oscillation]
\label{rem:omega}
\emph{[unconditional]} $M(x)=\Omega_{\pm}(x^{1/2-\eps})$ for every
$\eps>0$, because $\zeta$ has zeros with $\Re\rho=1/2$. Moreover the Mertens
conjecture is false: $\limsup M(x)x^{-1/2}>1$ and
$\liminf M(x)x^{-1/2}<-1$ (Odlyzko--te Riele~\cite{OtR}, with later
improvements). Hence $H_\nu\gg v_\nu^{1/2-\eps}$ for infinitely many
excursions, and neither $\Ap$ nor $\Am$ stabilizes.
\end{remark}

\subsection{Criteria equivalent to RH}
\label{subsec:criteria}

The equivalences below concern three increasingly local objects: the area of
$M$, the quadratic energy inside the excursion blocks, and the energy of the
black curve.  The analytic ingredient used for the converses is the $L^1$
Mellin criterion: if $\int_1^\infty|M(x)|x^{-\sigma-1}dx<\infty$ for every
$\sigma>1/2$, then $1/(s\zeta(s))=\int_1^\infty M(x)x^{-s-1}dx$ is analytic
in $\Re s>1/2$, so $\zeta$ has no zeros there.

\begin{proposition}[Excursionwise $L^1$ criterion]
\label{prop:L1}
\emph{[unconditional: an equivalence, RH is not assumed]} Let
$\beta_\nu=\sum_{n\in I_\nu}|M(n)|$. Then
\[
  \mathrm{RH}\iff \sum_{v_\nu\le x}\beta_\nu\ \ll_\eps\ x^{3/2+\eps}
  \quad\text{for every }\eps>0,
\]
the excursion in progress included. The sum does not depend on the signs
$\sigma_\nu$: RH is a statement about the profile $|M|$.
\end{proposition}

\begin{lemma}[Pruning of short excursions]
\label{lem:pruning}
\emph{[unconditional]} For every complete excursion,
\begin{equation}
\label{eq:trivial-beta}
  \beta_\nu\le H_\nu\ell_\nu\le \frac{\ell_\nu^2}{2}.
\end{equation}
Consequently, for every threshold $L_0\ge1$,
\begin{equation}
\label{eq:pruning}
  \sum_{\substack{v_\nu\le x\\ \ell_\nu\le L_0}}\beta_\nu
  \le \frac{L_0}{2}\sum_{\substack{v_\nu\le x\\ \ell_\nu\le L_0}}\ell_\nu
  \le \frac{xL_0}{2}.
\end{equation}
\end{lemma}

\begin{proof}
Since an excursion must climb from $0$ to $H_\nu$ and return to $0$,
$H_\nu\le\min(a_\nu,b_\nu)\le\ell_\nu/2$. Also
$|M(n)|\le H_\nu$ throughout the excursion, which gives
$\beta_\nu\le H_\nu\ell_\nu$. If $\ell_\nu\le L_0$, then
$\ell_\nu^2\le L_0\ell_\nu$; summing and using the disjointness of
complete excursions proves \eqref{eq:pruning}.
\end{proof}

\begin{corollary}[Long-excursion reduction]
\label{cor:long-reduction}
\emph{[unconditional: an equivalence, RH is not assumed]} Fix any
$\delta>0$ and put $L_0=x^{1/2-\delta}$. Then
\begin{equation}
\label{eq:long-reduction}
  \mathrm{RH}
  \iff
  \sum_{\substack{v_\nu\le x\\
                      \ell_\nu>x^{1/2-\delta}}}\beta_\nu
  +\beta_{\rm term}(x)
  \ll_\eps x^{3/2+\eps}
  \qquad\text{for every }\eps>0,
\end{equation}
where the terminal block is included only when it crosses the same length
threshold. In particular the number of complete excursions surviving the
cut is at most
\begin{equation}
\label{eq:number-long}
  \#\{\nu:v_\nu\le x,\ \ell_\nu>L_0\}\le \frac{x}{L_0}
  =x^{1/2+\delta}.
\end{equation}
More generally, the number with $L<\ell_\nu\le2L$ is at most $x/L$.
\end{corollary}

\begin{proof}
By Lemma~\ref{lem:pruning}, the deleted excursions contribute at most
$\tfrac12x^{3/2-\delta}$, which is below the $x^{3/2+\eps}$ budget in
Proposition~\ref{prop:L1}. Thus deleting them does not change the
criterion. The counting bounds follow from disjointness and
$\sum\ell_\nu\le x$.
\end{proof}

\begin{remark}[The square-root frontier]
\label{rem:sqrt-frontier}
Taking $L_0=\sqrt{x}$ in \eqref{eq:pruning} gives exactly
$x^{3/2}/2$. Thus $\sqrt{x}$ is the frontier at which the completely
trivial excursion bound already matches the RH-scale area budget. A cut
strictly below $\sqrt{x}$ makes the discarded contribution lower order.
Pushing the useful cut substantially above this frontier, or obtaining a
smaller contribution from the retained side by the same argument, requires
additional information such as a nontrivial height bound.
\end{remark}

\begin{heuristic}[Renewal-tail concentration of the area]
\label{heu:renewal-tail}
\emph{[heuristic]} Suppose the excursion-length tail has the renewal form
\begin{equation}
\label{eq:renewal-tail}
  \#\{\nu:v_\nu\le x,\ \ell_\nu>L\}
  \asymp c\sqrt{x}L^{-1/2},
\end{equation}
and suppose a typical excursion at scale $\ell\asymp L$ has
$\beta\asymp\sigma L^{3/2}$. Then a dyadic scale $[L,2L]$ contains
$\asymp c\sqrt{x}L^{-1/2}$ excursions and contributes
\begin{equation}
\label{eq:renewal-scale-area}
  \asymp c\sigma\sqrt{x}L
\end{equation}
to the total area. Hence the dyadic contributions grow linearly with the
length scale and are dominated by the longest excursions. Scales up to
$L$ contribute only order $L/x$ of the contribution at the terminal scale
$L\asymp x$.

For $N=455\,052\,511$, $\sqrt N\approx2.13\times10^4$ and
$Z(N)=87\,385$. Normalizing \eqref{eq:renewal-tail} by the observed
excursion count gives an order-of-magnitude prediction of roughly
$Z(N)N^{-1/4}\approx6.0\times10^2$ excursions longer than $\sqrt N$,
compared with all $87\,385$ excursions. This numerical count is a model
prediction, not a consequence of the pruning lemma.
\end{heuristic}

\begin{remark}[Spectral interpretation after pruning]
\label{rem:spectral-pruning}
The pruning lemma removes the short-excursion cloud without using any
information about the zeros of $\zeta$. What remains is precisely the
long-length regime in which a low-frequency approximation to the explicit
formula is potentially informative. In the spectral-resolution heuristic,
a truncation at ordinate $T$ resolves structures on length scales roughly
$\gtrsim x/T$. Thus the pruning step and the low-zero description act in
complementary regimes: the former removes a provably harmless short-scale
part, while the latter is relevant only to the retained long excursions.
This is a heuristic interpretation, not an unconditional transfer theorem
for excursion-selected intervals.
\end{remark}

\begin{remark}[What pruning does not solve]
\label{rem:pruning-limit}
The reduction changes where the difficulty lies, not its size on the
largest blocks. If $\ell_\nu\asymp x$, the trivial estimate
$\beta_\nu\le\ell_\nu^2/2$ is of order $x^2$, while the RH-scale
budget is $x^{3/2+\eps}$. A factor of essentially $x^{1/2}$ is still
missing, and there are only $O(1)$ such terminal-scale excursions, so no
averaging over many blocks is available. This is the regime in which the
local profile, the correlation moment $Q_\nu$, and the bridge geometry
must supply genuinely new information.
\end{remark}

\begin{theorem}[Aggregated internal-correlation criterion]
\label{thm:corr-criterion}
\emph{[unconditional: an equivalence, RH is not assumed]} With $Q_\nu$ and
$C_h^{(\nu)}$ as in \S\ref{subsec:internal-corr}, and with the terminal
block included through \eqref{eq:general-triangular},
\begin{equation}
\label{eq:RH-Q}
  \mathrm{RH}
  \iff
  \sum_{v_\nu\le x}Q_\nu+Q_{\rm term}(x)
  \ll_\eps x^{2+\eps}
  \qquad\text{for every }\eps>0.
\end{equation}
Equivalently,
\begin{equation}
\label{eq:RH-corr}
  \mathrm{RH}
  \iff
  -\sum_{v_\nu\le x}\ \sum_{h=1}^{\ell_\nu-1}
  hC_h^{(\nu)}+Q_{\rm term}(x)
  \ll_\eps x^{2+\eps}
  \qquad\text{for every }\eps>0.
\end{equation}
Thus RH is equivalent to an averaged bound for the first weighted moment of
the correlations internal to the excursions; no excursionwise bound
$Q_\nu\ll\ell_\nu^{2+\eps}$ is required.
\end{theorem}

\begin{proof}
By Proposition~\ref{prop:quadratic-decomp}, the left side of
\eqref{eq:RH-Q} is exactly $\sum_{n\le x}M(n)^2$.  Under RH,
$M(n)\ll_\eps n^{1/2+\eps}$, hence
$\sum_{n\le x}M(n)^2\ll_\eps x^{2+\eps}$ after renaming $\eps$.
Conversely, if the latter bound holds for every $\eps>0$, then
Cauchy--Schwarz gives
\[
  \sum_{n\le x}|M(n)|
  \le x^{1/2}\left(\sum_{n\le x}M(n)^2\right)^{1/2}
  \ll_\eps x^{3/2+\eps}.
\]
Partial summation gives the $L^1$ Mellin criterion for every
$\sigma>1/2$, and therefore RH.  Equation \eqref{eq:RH-corr} follows from
\eqref{eq:Qcorr} and the terminal-block identity.
\end{proof}

\begin{corollary}[A local sufficient condition]
\label{cor:local-corr-sufficient}
\emph{[unconditional implication]} If, for every $\eps>0$, every complete
excursion satisfies
\[
  Q_\nu=-\sum_{h=1}^{\ell_\nu-1}hC_h^{(\nu)}
  \ll_\eps \ell_\nu^{2+\eps},
\]
and the terminal block satisfies the analogous bound, then RH holds.
Equivalently, it is enough to have
\[
  \sum_{h<\ell_\nu}h\bigl(2D_h^{(\nu)}-P_h^{(\nu)}\bigr)
  \ll_\eps \ell_\nu^{2+\eps}
\]
for every excursion.
\end{corollary}

\begin{remark}[Calibration: a Chowla-scale condition]
\label{rem:chowla-scale}
The local condition is sufficient but not weak.  If the correlations
$C_h^{(\nu)}$ had a typical size $B$ with essentially uncorrelated signs,
then
\[
  \left|\sum_{h<L}hC_h^{(\nu)}\right|
  \asymp B\left(\sum_{h<L}h^2\right)^{1/2}
  \asymp B L^{3/2}.
\]
The target $O(L^2)$ therefore corresponds, in such a generic regime, to
$B=O(\sqrt L)$: square-root cancellation in the pair correlations averaged
over the lag variable.  Its advantage over pointwise Chowla is not a lower
cancellation scale, but the freedom to tolerate individual bad lags if they
compensate in the weighted aggregate.

At the unconditioned density level one expects
$\sum_{h<L}hP_h$ to be of order $L^3$ (the crude independent-squarefree
scale gives $\tfrac12(6/\pi^2)^2L^3\approx0.18L^3$).  Hence
\eqref{eq:Q-P-D} asks for cancellation of the two order-$L^3$ contributions
down to order $L^2$.  This is a calibration of the required strength, not a
claim that Mirsky estimates on fixed intervals apply directly to the
excursion-selected intervals.
\end{remark}

\begin{remark}[The conditioning obstruction and the missing third constraint]
\label{rem:conditioning}
The interval $I_\nu=(u_\nu,v_\nu]$ is not fixed independently of $\mu$:
its endpoints are selected by the condition that the partial sum vanishes,
and, after orientation, the whole block obeys
\[
  S_t>0\quad(1\le t<\ell_\nu),\qquad S_{\ell_\nu}=0.
\]
Thus estimates for $\sum_j\mu(j)\mu(j+h)$ on fixed intervals, or even on
average over fixed intervals, cannot be transferred to $I_\nu$ without
controlling this selection bias.  The conditioning itself forces
\[
  -\sum_hhC_h^{(\nu)}=Q_\nu\ge \frac{\beta_\nu^2}{\ell_\nu},
\]
so a local correlation-energy bound implies the local area bound, but not
conversely unless the profile is known to be sufficiently flat.

Equations \eqref{eq:zeroth-corr} and \eqref{eq:Qcorr} give two exact
constraints on the zeroth and first weighted moments of the internal
correlations.  The remaining structural problem is to find a third
constraint that prevents the negative correlation mass from moving too far
towards lags $h\asymp\ell_\nu$.  A natural place to seek it is in the bridge
geometry---for example through the climb/descent asymmetry and through a
defect parameter of the form $\delta=\ell/2-H$ (or its version using the
number of non-zero steps).  Large transport of negative mass to long lags
corresponds precisely to a more ballistic excursion profile.
\end{remark}

\begin{theorem}[Dirichlet energy of the black curve]
\label{thm:energy}
\emph{[unconditional: an equivalence, RH is not assumed]} Let $\mathcal{J}$
be the linear interpolation of $\Anet$ between consecutive critical points
and
\[
  E(x)=\int_0^x(\mathcal{J}')^2
  =\sum_{v_\nu\le x}\frac{\bigl(\Anet(v_\nu)-\Anet(u_\nu)\bigr)^2}{\ell_\nu}
  +\frac{(x-u)\max_{u<n\le x}|M(n)|}{4},
\]
where $(u,x]$ is the leg in progress. By Lemma~\ref{lem:jump},
$\bigl(\Anet(v_\nu)-\Anet(u_\nu)\bigr)^2/\ell_\nu=H_\nu\ell_\nu/4$, so that
$E(x)=\tfrac14\bigl(\sum_{v_\nu\le x}H_\nu\ell_\nu+(x-u)H_{\mathrm{run}}\bigr)$.
Then
\[
  \sum_{n\le x}|M(n)|\ \le\ 4E(x)\ \le\ x\max_{n\le x}|M(n)|,
\]
and
\[
  \mathrm{RH}\iff E(x)\ll_\eps x^{3/2+\eps}\quad\text{for every }\eps>0 .
\]
\end{theorem}

\begin{proof}
The lower bound is $|M(n)|\le H_\nu$ summed over each excursion, and the
upper one is $H_\nu\le\max_{n\le x}|M(n)|$ together with
$\sum\ell_\nu\le x$. If RH holds, $\max_{n\le x}|M(n)|\ll_\eps x^{1/2+\eps}$
and the upper bound gives $E(x)\ll x^{3/2+\eps}$. Conversely, if
$E(x)\ll x^{3/2+\eps}$, the lower bound gives
$\sum_{n\le x}|M(n)|\ll x^{3/2+\eps}$; by partial summation,
$\int_1^\infty|M(x)|x^{-\sigma-1}dx<\infty$ for $\sigma>1/2$, and the Mellin
criterion concludes.
\end{proof}

\begin{remark}
\label{rem:average}
$E$ is an average: $4E(x)=\sum H_\nu\ell_\nu$ is the excursion height
weighted by $\ell_\nu$. It allows isolated excursions with
$H_\nu\gg v_\nu^{1/2}$ provided they are short or rare, and only forbids the
weighted mass from exceeding $x^{3/2+\eps}$. It is moreover a sum of
nonnegative local terms, with no sign cancellation.
\end{remark}


\subsection{Prime-information horizon of an excursion}
\label{subsec:prime-horizon}

This subsection isolates an elementary but useful information boundary.  It
separates what is forced by the prime data already available at the beginning
of an excursion from what would genuinely require new prime information after
the excursion has started.

\begin{lemma}[Square horizon from the primes known at the start]
\label{lem:square-horizon}
\emph{[unconditional]} Let $u\ge2$.  Knowledge of the complete set of primes
$p\le u$ determines $\mu(n)$ for every integer $n\le u^2$.
\end{lemma}

\begin{proof}
Remove from $n$ all prime factors $p\le u$, keeping track of their
multiplicities.  If a known prime occurs twice then $\mu(n)=0$.  Otherwise the
remaining factor $q$ is either $1$ or a single prime $>u$: it cannot be a
product of two integers $>u$, because then $q>u^2\ge n$.  Thus the parity of
the number of distinct prime factors, and hence $\mu(n)$, is determined using
no prime information beyond $u$.
\end{proof}

\begin{definition}[Information-self-contained and information-expanding excursions]
\label{def:information-excursion}
For a complete excursion $I_\nu=(u_\nu,v_\nu]$, call $I_\nu$
\emph{information-self-contained at its start} if
\[
  v_\nu\le u_\nu^2,
\]
and \emph{information-expanding} if
\[
  v_\nu>u_\nu^2.
\]
Equivalently, writing $L_\nu=v_\nu-u_\nu$, the exceptional condition is
$L_\nu>u_\nu^2-u_\nu$.  By Lemma~\ref{lem:square-horizon}, every
information-self-contained excursion is completely determined by the primes
already known at its starting point $u_\nu$; no prime discovered after
$u_\nu$ is needed as new input to determine the values of $\mu$ throughout
that excursion.
\end{definition}

For a general excursion, the fraction of its integer time interval that is
unconditionally determined from the prime data available at its start is at
least
\begin{equation}
\label{eq:known-fraction}
  \vartheta_\nu
  =\min\!\left\{1,
  \frac{u_\nu^2-u_\nu}{v_\nu-u_\nu}\right\}.
\end{equation}
Thus the only excursions for which the known-prime horizon fails to cover the
whole block are precisely the information-expanding excursions.

\begin{proposition}[Double-exponential separation of information-expanding excursions]
\label{prop:double-exp-separation}
\emph{[unconditional]} Let
\[
  a_1<a_2<a_3<\cdots
\]
be the starting points, restricted to $a_j\ge2$, of the information-expanding
excursions.  Then
\begin{equation}
\label{eq:double-exp-separation}
  a_{j+1}>a_j^2.
\end{equation}
Consequently
\begin{equation}
\label{eq:double-exp-growth}
  a_j>a_1^{\,2^{j-1}},
\end{equation}
and the number
\[
  N_{\rm new}(X)
  :=\#\{\nu:u_\nu\le X,\ v_\nu>u_\nu^2\}
\]
satisfies
\begin{equation}
\label{eq:loglog-exception-count}
  N_{\rm new}(X)=O(\log\log X).
\end{equation}
\end{proposition}

\begin{proof}
If an information-expanding excursion starts at $a_j$, it ends at a point
$v>a_j^2$.  Every later excursion starts at that endpoint or later.  Hence the
next information-expanding excursion, if one exists, has
$a_{j+1}>a_j^2$.  Iteration gives \eqref{eq:double-exp-growth}, and solving
$a_1^{2^{j-1}}\le X$ for $j$ gives \eqref{eq:loglog-exception-count}.
\end{proof}

\begin{remark}[What the unconditional separation does and does not imply]
\label{rem:unconditional-information-separation}
The bound $N_{\rm new}(X)=O(\log\log X)$ is unconditional, but it does not imply
that only finitely many information-expanding excursions exist.  An infinite
sequence is still logically possible, provided its starting points satisfy the
double-exponential separation $a_{j+1}>a_j^2$.  No eventual-zero statement is
used below.
\end{remark}

\subsection{Propagation of prime information beyond the square horizon}
\label{subsec:prime-copy-propagation}

The square horizon identifies where information not already certified at the
starting point can first become relevant.  The next question is quantitative:
how much can genuinely new prime information change the height of an excursion,
once it is combined with the part of the trajectory already determined at the
start?  In this subsection it is convenient to write
\[
  m(x):=M(x)=\sum_{n\le x}\mu(n)
\]
for the scalar Mertens sum.

\begin{definition}[First-order new-prime channel]
\label{def:first-order-channel}
Fix a starting point $u\ge2$.  For a prime $p>u^2$ and a point $t<p^2$, define
\[
  \Delta_p(t):=\sum_{a\le t/p}\mu(ap).
\]
Since $a<t/p<p$, one has $(a,p)=1$, and therefore
\begin{equation}
\label{eq:single-prime-copy}
  \boxed{\qquad
  \Delta_p(t)=-m\!\left(\left\lfloor\frac{t}{p}\right\rfloor\right).
  \qquad}
\end{equation}
Thus a new prime does not contribute an arbitrary string of equal signs: it
reproduces, with reversed sign, a prefix of the Mertens walk that is already
known at the smaller scale $t/p$.
\end{definition}

\begin{proposition}[Forced $+1,-1,0$ content of one prime cone]
\label{prop:prime-cone-signs}
\emph{[unconditional]} Put $y=\lfloor t/p\rfloor<p$ and
\[
  Q(y):=\sum_{a\le y}\mu(a)^2.
\]
Among the values $\mu(ap)$, $1\le a\le y$, the exact numbers of $+1,-1$ and
$0$ are
\begin{equation}
\label{eq:prime-cone-counts}
  N_+(y)=\frac{Q(y)-m(y)}2,\qquad
  N_-(y)=\frac{Q(y)+m(y)}2,\qquad
  N_0(y)=y-Q(y).
\end{equation}
Consequently
\[
  N_+(y)-N_-(y)=-m(y)=\Delta_p(t).
\]
In particular, using the squarefree density and the prime number theorem,
\[
  \frac{N_+(y)}y\to\frac3{\pi^2},\qquad
  \frac{N_-(y)}y\to\frac3{\pi^2},\qquad
  \frac{N_0(y)}y\to1-\frac6{\pi^2},
\]
and
\begin{equation}
\label{eq:relative-prime-cone-bias}
  \frac{|\Delta_p(t)|}{y}=\frac{|m(y)|}{y}\longrightarrow0.
\end{equation}
The cancellation in \eqref{eq:relative-prime-cone-bias} is only relative to
the full cone size $y$; it is not a square-root bound.
\end{proposition}

\begin{proof}
For $a<p$, multiplicativity gives $\mu(ap)=-\mu(a)$.  The identities in
\eqref{eq:prime-cone-counts} follow from
$N_+(\mu)+N_-(\mu)=Q(y)$ and
$N_+(\mu)-N_-(\mu)=m(y)$ before the sign reversal.  The asymptotic
proportions use $Q(y)=6y/\pi^2+O(\sqrt y)$ and $m(y)=o(y)$.
\end{proof}

\begin{proposition}[Linear first-order regime]
\label{prop:linear-new-primes}
\emph{[unconditional]} Suppose
\begin{equation}
\label{eq:linear-window}
  u^2<t<u^4.
\end{equation}
Then two distinct primes $p,q>u^2$ cannot divide the same integer $n\le t$,
and no square $p^2$ with $p>u^2$ can occur below $t$.  Hence the first-order
new-prime channels are disjoint and their total contribution is exactly
\begin{equation}
\label{eq:first-order-total}
  \boxed{\qquad
  E_u^{(1)}(t)
  =-\sum_{u^2<p\le t}
  m\!\left(\left\lfloor\frac{t}{p}\right\rfloor\right).
  \qquad}
\end{equation}
This formula concerns the first-order channels generated by primes $p>u^2$;
it is not asserted here to be a complete parametrization of every possible
residual-information mechanism beyond the square horizon.
\end{proposition}

\begin{proof}
If $p,q>u^2$, then $pq>u^4>t$, and likewise $p^2>u^4>t$.
Thus the sets of multiples of the different $p$ occurring below $t$ are
disjoint.  Equation \eqref{eq:first-order-total} follows by summing
\eqref{eq:single-prime-copy}.
\end{proof}

\begin{definition}[Copy length and multiplicity]
\label{def:copy-multiplicity}
Set
\[
  R=R(t,u):=\frac{t}{u^2}.
\]
For an integer $1\le y\le\lfloor R\rfloor$, let
\[
  N_y(t,u)
  :=\#\left\{p>u^2:\ p\le t,\
       \left\lfloor\frac{t}{p}\right\rfloor=y\right\}.
\]
Equivalently,
\begin{equation}
\label{eq:Ny-exact}
  N_y(t,u)
  =\#\left\{p:\
      \max\!\left(u^2,\frac{t}{y+1}\right)<p\le\frac{t}{y}
    \right\}.
\end{equation}
Therefore the first-order contribution has the exact grouped form
\begin{equation}
\label{eq:first-order-grouped}
  \boxed{\qquad
  E_u^{(1)}(t)=-\sum_{1\le y\le R}N_y(t,u)m(y).
  \qquad}
\end{equation}
The total number of copies is
\begin{equation}
\label{eq:copy-total}
  K(t,u)=\sum_yN_y(t,u)=\pi(t)-\pi(u^2).
\end{equation}
\end{definition}

\begin{remark}[How many copies and how long they are]
\label{rem:copy-count-asymptotic}
The following are unconditional PNT asymptotics, not probabilistic
assumptions.  If $t=Ru^2$ with fixed $R>1$ and $u\to\infty$, then
\begin{equation}
\label{eq:K-fixed-R}
  K(t,u)\sim \frac{(R-1)u^2}{2\log u}.
\end{equation}
For each fixed $y<R$,
\begin{equation}
\label{eq:Ny-fixed-R}
  N_y(t,u)
  \sim
  \frac{Ru^2}{2\log u}\frac1{y(y+1)}.
\end{equation}
No uniformity in growing $R$ is claimed in \eqref{eq:Ny-fixed-R}.  The exact
formula to use outside the fixed-$R$ regime is \eqref{eq:Ny-exact}.
\end{remark}

\begin{remark}[The first copy-length transitions]
\label{rem:first-copy-transitions}
The parameter $R=t/u^2$ measures how far the excursion has travelled beyond
the square information horizon.  If $1<R<2$, every new-prime copy has length
$1$, so $m(1)=1$ and
\[
  E_u^{(1)}(t)=-\bigl(\pi(t)-\pi(u^2)\bigr).
\]
If $2<R<3$, copies of length $2$ appear, but $m(2)=0$, so that channel is
exactly neutral.  At $3<R<4$, copies of length $3$ enter and
$m(3)=-1$, creating channels of the opposite net sign.  Thus the forced
$+1,-1,0$ structure begins to screen the first-order contribution as the
available copy lengths diversify.  This statement is exact at each fixed
$t$ through \eqref{eq:first-order-grouped}; it does not assert an RH-scale
bound.
\end{remark}

\begin{lemma}[Ideal copy-weight cancellation coefficient]
\label{lem:copy-weight-coefficient}
\emph{[unconditional]} Define
\begin{equation}
\label{eq:AR-def}
  A(R):=\sum_{1\le y\le R}\frac{m(y)}{y(y+1)}.
\end{equation}
Then
\begin{equation}
\label{eq:AR-exact}
  \boxed{\qquad
  A(R)=\sum_{n\le R}\frac{\mu(n)}n-\frac{m(R)}{R+1}.
  \qquad}
\end{equation}
Moreover
\begin{equation}
\label{eq:AR-tail}
  A(R)=-\sum_{n>R}\frac{m(n)}{n(n+1)},
\end{equation}
where the convergent identity uses
$\sum_{n\ge1}\mu(n)/n=0$, equivalently the prime number theorem.  Hence
\begin{equation}
\label{eq:AR-zero}
  A(R)\longrightarrow0.
\end{equation}
\end{lemma}

\begin{proof}
Reverse the order of summation:
\[
  A(R)
  =\sum_{n\le R}\mu(n)
    \sum_{y=n}^{R}\frac1{y(y+1)}
  =\sum_{n\le R}\mu(n)
    \left(\frac1n-\frac1{R+1}\right),
\]
which gives \eqref{eq:AR-exact}.  Letting $R\to\infty$ and using the prime
number theorem gives the infinite sum $0$; subtracting the finite partial
sum gives \eqref{eq:AR-tail}.
\end{proof}

\begin{corollary}[Unconditional decay available for the ideal coefficient]
\label{cor:AR-VK}
Using the Vinogradov--Korobov zero-free-region bound
\[
  m(x)\ll
  x\exp\!\left[-c(\log x)^{3/5}(\log\log x)^{-1/5}\right],
\]
there is $c_1>0$ such that
\begin{equation}
\label{eq:AR-VK}
  A(R)\ll
  \exp\!\left[-c_1(\log R)^{3/5}(\log\log R)^{-1/5}\right].
\end{equation}
Thus the idealized coherent fraction of increasingly long copies tends to
zero unconditionally, but only at a sub-power rate.
\end{corollary}

\begin{remark}[Main-term interpretation and its limitation]
\label{rem:AR-main-term}
For fixed $R$, substituting the PNT approximation
\eqref{eq:Ny-fixed-R} into the exact grouped sum
\eqref{eq:first-order-grouped} gives the first-order main term
\begin{equation}
\label{eq:E-main-fixed-R}
  E_u^{(1)}(Ru^2)
  \sim -\frac{Ru^2}{2\log u}\,A(R)
  \qquad(u\to\infty,\ R\ \text{fixed}).
\end{equation}
The coefficient $A(R)$ therefore measures the signed cancellation forced by
the ensemble of copy lengths in this fixed-$R$ asymptotic.  When $R$ grows
with $u$, \eqref{eq:first-order-grouped} remains exact but
\eqref{eq:E-main-fixed-R} must not be used without a uniform prime-counting
estimate for the weights $N_y(t,u)$.
\end{remark}

\begin{definition}[Remaining margin to a dangerous height]
\label{def:remaining-margin}
Let $B_u(t)$ denote any baseline trajectory contribution already fixed by the
information retained at the start $u$, before adding the first-order channels
\eqref{eq:first-order-total}.  For a target
$T_\eps(t)=t^{1/2+\eps}$ define the remaining margin
\begin{equation}
\label{eq:margin-def}
  G_{u,\eps}(t):=
  \max\bigl\{0,T_\eps(t)-|B_u(t)|\bigr\}.
\end{equation}
The first-order new-prime information can be causally responsible for crossing
the target at $t$ only if its signed contribution has magnitude at least this
margin.  A sign-free sufficient test excluding such responsibility is
\begin{equation}
\label{eq:causal-capacity}
  C_u^{(1)}(t):=
  \sum_{y\le R}N_y(t,u)|m(y)|<G_{u,\eps}(t),
\end{equation}
whereas the exact first-order signed test is
\begin{equation}
\label{eq:causal-signed}
  \left|\sum_{y\le R}N_y(t,u)m(y)\right|
  <G_{u,\eps}(t).
\end{equation}
If \eqref{eq:causal-signed} holds, the first-order channels cannot supply the
missing height at that point.  No claim is made that the baseline/new-channel
decomposition by itself controls all higher-order residual information.
\end{definition}

\begin{remark}[Why long-copy cancellation does not prove RH]
\label{rem:copy-cancellation-not-RH}
The limit $A(R)\to0$ is only a statement of \emph{relative} cancellation.
An $o(1)$ coherent fraction multiplying a quantity of order $t/\log t$ can
still be vastly larger than $t^{1/2+\eps}$.  To see the required scale,
formally compare the ideal first-order size
\[
  \frac{t}{\log t}|A(R)|
\]
with $t^{1/2+\eps}$.  A sufficient inequality would be
\begin{equation}
\label{eq:AR-required-t}
  |A(R)|\ll t^{-1/2+\eps}\log t.
\end{equation}
The unconditional decay \eqref{eq:AR-VK} is much weaker than any fixed
negative power of $R$.  Thus the known cancellation does not bound the
``deterministic randomness'' at the RH scale.
\end{remark}

\begin{remark}[Power-law calibration of the missing decay]
\label{rem:power-calibration-copy}
Write $t=u^\alpha$ with $\alpha>2$.  Then
$R=t/u^2=u^{\alpha-2}$.  The formal requirement
\eqref{eq:AR-required-t} becomes, up to logarithmic factors,
\begin{equation}
\label{eq:AR-required-R}
  |A(R)|\lesssim R^{-\delta(\alpha,\eps)},
  \qquad
  \delta(\alpha,\eps)
  =\frac{\alpha(1/2-\eps)}{\alpha-2}.
\end{equation}
For $\eps$ near $0$ this asks, respectively, for exponents $3/2$ at
$t=u^3$, $1$ at $t=u^4$, $3/4$ at $t=u^6$, and $5/8$ at $t=u^{10}$;
as $\alpha\to\infty$ the required exponent decreases only to $1/2$.
Thus increasing the excursion length relaxes the needed cancellation, but
never below the square-root barrier.  This is a calibration of the ideal
first-order mechanism, not an independent proof of an RH-scale bound.
\end{remark}

\begin{remark}[Hierarchy of interaction scales]
\label{rem:new-prime-hierarchy}
The square horizon $u^2$ is the first scale at which first-order prime
information outside the starting data can enter.  Two primes $>u^2$ cannot
occur together below $u^4$; three cannot occur together below $u^6$, and in
general a product of $k$ such primes lies above $u^{2k}$.  Hence
\[
  u^2,\ u^4,\ u^6,\ldots
\]
form natural thresholds for successive orders of interaction among these
new-prime channels.  The interval $u^2<t<u^4$ is distinguished because the
first-order superposition \eqref{eq:first-order-total} is exactly linear.
\end{remark}


\subsection{Model of the black curve}
\label{subsec:model}

\begin{heuristic}[Shape of $\Anet$]
\label{heu:model}
\emph{[heuristic; assumes RH, simple zeros, convergence of the
Fourier--Bohr expansion and a constant fill factor]} Writing
$M(t)=t^{1/2}F(\log t)$ with
$F(u)=2\Re\sum_{\gamma>0}e^{i\gamma u}/(\rho\zeta'(\rho))$, and
approximating $\Ms\approx\kappa\,\sgn(M)|M|^{1/2}$ with $\kappa$ constant,
one gets $\Ms(t)\approx\kappa\,t^{1/4}G(\log t)$ with
$G=\sgn(F)|F|^{1/2}$. If $G(u)=\sum_\lambda g_\lambda e^{i\lambda u}$, then
\[
  \Anet(x)\approx\kappa\,x^{5/4}\sum_\lambda \frac{g_\lambda\,x^{i\lambda}}{5/4+i\lambda}.
\]
This reproduces the exponent $5/4$, the smoothness of the waves (the factor
$|5/4+i\lambda|^{-1}$ damps the high frequencies) and the quasi-periodicity
in $\log x$: with dominant frequency $\gamma_1=14.1347$, the period is
$2\pi/\gamma_1$ in $\log x$, a factor $\approx1.56$ in $x$.
\end{heuristic}

\subsection{Numerical data}
\label{subsec:data}

With $N=455\,052\,511$ and $Z(N)=87\,385$ excursions:
\[
  \Ap(N)=6.3449\cdot10^{9},\qquad
  \Am(N)=-7.0877\cdot10^{9},\qquad
  \Anet(N)=-7.4279\cdot10^{8},
\]
\[
  \max_\nu \frac{H_\nu}{\sqrt{\ell_\nu}}=2.0743,\qquad
  \max_\nu \frac{\beta_\nu}{\ell_\nu^{3/2}}=1.0214 .
\]
For the pruned program the more relevant diagnostic is the same maximum
restricted to $\ell_\nu>\sqrt N$ (or, more generally, to
$\ell_\nu>N^{1/2-\delta}$). The value of that filtered maximum is not
reported here and should be recomputed from the excursion census rather
than inferred from the unfiltered maximum.
The second normalization is the one singled out by
\eqref{eq:beta-Q}: a local $Q_\nu=O(\ell_\nu^2)$ law gives exactly
$\beta_\nu=O(\ell_\nu^{3/2})$, whereas the same quadratic-energy control
only yields the weaker height law $H_\nu=O(\ell_\nu^{2/3})$.
The bound of Lemma~\ref{lem:transfer} with
$c=\max_{10^3\le n\le N}|M(n)|n^{-1/2}$ lies about $30$ times above
$|\Anet(N)|$. In this range $|M(x)|\lesssim 0.5\sqrt{x}$, well inside any
bound at the RH scale, so the data do not separate $x^{1/2}$ from
$x^{1/2+\eps}$: the factor $x^{\eps}$ with $\eps=0.05$ equals
$\approx 2.7$ here.

\begin{thebibliography}{9}
\bibitem{Titchmarsh} E.~C. Titchmarsh, \emph{The Theory of the Riemann
Zeta-function}, 2nd ed., revised by D.~R. Heath-Brown, Oxford Univ. Press,
1986, \S14.27.
\bibitem{OtR} A.~M. Odlyzko and H.~J.~J. te Riele, \emph{Disproof of the
Mertens conjecture}, J. Reine Angew. Math. \textbf{357} (1985), 138--160.
\bibitem{Ng} N.~Ng, \emph{The distribution of the summatory function of the
Möbius function}, Proc. London Math. Soc. (3) \textbf{89} (2004), 361--389.
\end{thebibliography}

\end{document}
